%% kapa: exam=1 round=30 year=2019 period=2 form=A subject=appraisal_law subject_ko=감정평가관계법규 seq=37 number=37
%% source: Q-Net 감정평가사 기출문제 (cst00309) artlSeq=5205205 file=extracted/1008/5205205/2019년 제30회 감정평가사 1차시험 문제지/2019년 제30회 감정평가사 1차시험 문제지/2019년도 제30회 감정평가사 1차 2교시 A형.hwp
%% conversion: kordoc 4.24.51 -> tools/gen1.py (auto); review=proofread-sonnet; flags=none
\begin{question}{37}
도시 및 주거환경정비법령상 조합임원에 관한 설명으로 옳은 것은?
\begin{choices}
  \item 조합임원이 금고 이상의 형의 집행유예를 받고 그 유예기간 중에 있는 경우에는 총회의 의결을 거쳐 해임된다.
  \item 조합임원은 조합원 10분의 1 이상의 요구로 소집된 총회에서 조합원 과반수의 출석과 출석 조합원 과반수의 동의를 받아 해임할 수 있다.
  \item 조합장 또는 이사가 자기를 위하여 조합과 계약이나 소송을 할 때에는 대의원회의 의장이 조합을 대표한다.
  \item 조합임원의 임기는 정관으로 정하되, 연임할 수 없다.
  \item 조합의 정관에는 조합임원 업무의 분담 및 대행 등에 관한 사항은 포함되지 아니한다.
\end{choices}
\end{question}
\answer{2}
